\section{Conclusion}

\noindent Cet article a présenté un cadre complet d'apprentissage profond pour la
détection automatique des fuites d'eau dans les réseaux urbains de distribution,
à partir de données capteurs en séries temporelles multivariées. L'approche
proposée intègre une ingénierie de caractéristiques guidée par le domaine, une
construction de séquences par fenêtre glissante, et une architecture hybride
CNN-BiLSTM-Attention afin de capturer efficacement les dépendances temporelles
locales et de long terme.

\vspace{0.2cm}

\noindent Les résultats expérimentaux montrent que le modèle atteint de fortes
performances, avec un compromis favorable entre précision, rappel et taux de
fausses alertes. Les comparaisons de benchmark confirment également que le
modèle proposé offre de meilleures performances opérationnelles que les
références classiques et de deep learning, ce qui met en évidence l'efficacité
de l'architecture hybride.

\vspace{0.2cm}

\noindent Les résultats soulignent l'importance de combiner des techniques de
modélisation avancées avec un prétraitement robuste des données, une évaluation
sans fuite de données et une calibration du seuil afin de développer des
systèmes de surveillance fiables et déployables.

\vspace{0.2cm}

\noindent Les travaux futurs porteront sur l'évaluation du modèle sur des jeux de
données plus diversifiés, l'extension du benchmark à des approches
supplémentaires, et l'amélioration de l'interprétabilité. Par ailleurs,
l'intégration de la localisation des fuites et de l'estimation de leur gravité
constitue une piste prometteuse pour renforcer l'applicabilité pratique.

\vspace{0.2cm}

\noindent Globalement, ce travail contribue au développement de solutions
intelligentes et fiables pour la surveillance urbaine de l'eau, en soutenant une
meilleure gestion des infrastructures et la réduction des pertes en eau.

\vspace{0.2cm}

\noindent \textbf{Applications potentielles.}
Le cadre proposé peut être intégré dans des systèmes de gestion intelligente de
l'eau basés sur l'Internet des objets (IoT), pour la détection et la
surveillance des fuites en temps réel. En traitant en continu les flux de
données capteurs, le modèle peut fournir des alertes précoces d'événements de
fuite, permettant une réaction plus rapide et une réduction des pertes d'eau.

\vspace{0.2cm}

\noindent En outre, le système peut être étendu pour supporter des mécanismes de
contrôle automatisés, tels que la gestion dynamique des vannes. En cas de fuite
détectée, le modèle peut déclencher la fermeture automatique des vannes ou la
régulation du débit afin d'isoler les segments affectés du réseau, et ainsi
minimiser les dommages et les perturbations opérationnelles.

\vspace{0.2cm}

\noindent Une telle intégration permet de passer d'une maintenance réactive à une
gestion proactive et autonome des infrastructures, en améliorant la fiabilité du
système et l'efficacité opérationnelle. L'approche proposée constitue ainsi un
brique essentielle pour les systèmes intelligents de distribution d'eau de
nouvelle génération.